\documentclass{article}
\usepackage{amsmath}
\usepackage{amssymb}
\usepackage{indentfirst}
\usepackage{graphicx}
\usepackage{color}
\usepackage{fancyhdr}
\usepackage{epstopdf}
\usepackage{indentfirst}
\usepackage{geometry}
\geometry{left=2.5cm,right=2.5cm,top=2.5cm,bottom=2.5cm}

\title{8.511 Problem Set 1}
\author{Yijun Jiang}
%\email{yjjiang@mit.edu}
\date{\today}

\pagestyle{fancy}
\lhead{Yijun Jiang}
\rhead{8.511 Problem Set 1}

\begin{document}
\maketitle
\section{Constructive Interference and Bravais Lattice}
\subsection{Part (a)}
Since $e^{i\vec{Q}\cdot\vec{R}_j}=e^{i\alpha}$, $\vec{Q}\cdot\vec{R}_j$ only differs from $\alpha$ by an integer multiple of $2\pi$. The same is true if we replace $\vec{R}_j$ with $2\vec{R}_j$, for $2\vec{R}_j$ is also a lattice point. Therefore, we have
\begin{align*}
	&\vec{Q}\cdot\vec{R}_j=\alpha+2k\pi\\
	&\vec{Q}\cdot2\vec{R}_j=\alpha+2k'\pi
\end{align*}
where $k,k'\in\mathbb{Z}$. Subtracting from the second equation twice the first one, we get
\begin{equation*}
	\alpha=2(k'-2k)\pi
\end{equation*}
Therefore, $\alpha$ is an integer multiple of $2\pi$. Since $\alpha$ is a phase and can be restricted to the interval $[0,2\pi)$, it is clear that constructive interference requires $\alpha=0$.

\subsection{Part (b)}
Let $\vec{Q}=\vec{Q}_0$ be an arbitrary solution to the original equation. Since $\vec{b}_1,\vec{b}_2,\vec{b}_3$ span the reciprocal space, there are $x_1,x_2,x_3\in\mathbb{R}$ such that
\begin{equation*}
	\vec{Q}_0=x_1\vec{b}_1+x_2\vec{b}_2+x_3\vec{b}_3
\end{equation*}
To prove the uniqueness of the general solution $\vec{Q}\in\{\vec{G}|\vec{G}=m_1\vec{b}_1+m_2\vec{b}_2+m_3\vec{b}_3,m_1,m_2,m_3\in\mathbb{Z}\}$, it suffices to show that $x_1,x_2,x_3\in\mathbb{Z}$.

Since $e^{i\vec{Q}_0\cdot\vec{a}_1}=e^{i2\pi x_1}=1$, $x_1\in\mathbb{Z}$. Similarly taking $\vec{R}_j=\vec{a}_2,\vec{a}_3$, we have $x_2,x_3\in\mathbb{Z}$. This proves the uniqueness of the solution.

\section{Problem 2}
\subsection{Part (a)}
Proof: The volume of the Bravais lattice is $V=\vec{a}_1\cdot(\vec{a}_2\times\vec{a}_3)$. The volume of the first Brillouin zone is
\begin{align*}
	V_k&=\vec{b}_1\cdot(\vec{b}_2\times\vec{b}_3)\\
	&=(2\pi)^3\frac{\vec{a}_2\times\vec{a}_3}{\vec{a}_1\cdot(\vec{a}_2\times\vec{a}_3)}\frac{\vec{a}_3\times\vec{a}_1}{\vec{a}_2\cdot(\vec{a}_3\times\vec{a}_1)}\frac{\vec{a}_1\times\vec{a}_2}{\vec{a}_3\cdot(\vec{a}_1\times\vec{a}_2)}\\
	&=\frac{(2\pi)^3}{V^3}(\vec{a}_2\times\vec{a}_3)\cdot(((\vec{a}_3\times\vec{a}_1)\cdot\vec{a}_2)\vec{a}_1-((\vec{a}_3\times\vec{a}_1)\cdot\vec{a}_1)\vec{a}_2)\\
	&=\frac{(2\pi)^3}{V^3}(\vec{a}_1\cdot(\vec{a}_2\times\vec{a}_3))(\vec{a}_2\cdot(\vec{a}_3\times\vec{a}_1))\\
	&=\frac{(2\pi)^3}{V}
\end{align*}

Therefore, $V_kV=(2\pi)^3$.

\subsection{Part (b)}
A conventional unit cell of the fcc lattice contains 4 atoms, so $4V^{fcc}=a^3$. On the other hand, a conventional unit cell of the bcc lattice contains 2 atoms, so $2V^{bcc}=a^3$.

The reciprocal lattice of fcc is bcc. Let its linear dimension be $b$. Then the first Brillouin zone has volume $V_k^{fcc}=b^3/2$. From the previous part, we have
\begin{equation*}
	\left(\frac{1}{2}b^3\right)\left(\frac{1}{4}a^3\right)=(2\pi)^3
\end{equation*}
whose solution gives $b=\frac{4\pi}{a}$.

\section{Problem 3}
\subsection{Part (a)}
The Bravais lattice and the corresponding unit cell is sketched below. The unit cell has a squared base whose linear scale is $a$. The primitive vectors are $\vec{a}_1=(a,0,0),\vec{a}_2=(0,a,0),\vec{a}_3=(0,0,c)$. The basis contains three vectors: $\vec{d}_1^{Cu}=(0,0,0),\vec{d}_1^{O}=(a/2,0,0),\vec{d}_2^{O}=(0,a/2,0)$.
\begin{figure}[!htbp]
	\centering
	\includegraphics[width=12cm]{blank.png}
\end{figure}

\subsection{Part (b)}
The primitive cell of the up-and-down CuO$_2$ layer is sketched below. Its area is $2a^2$. So the lattice spacing on the CuO$_{2}$ plane is $\sqrt{2}a$. The primitive vectors are (since I am only talking about one layer, the third dimension is left out) $\vec{a}_1=(a,-a)$ and $\vec{a}_2=(a,a)$.

The reciprocal lattice is also a square lattice. $\vec{b}_1=(\pi/a,-\pi/a)$ and $\vec{b}_2=(\pi/a,\pi/a)$.

Notice that the flat CuO$_2$ plane has a reciprocal lattice that is twice as large. This means, when the distortion vanishes, half of the $G$ points are ``lost'', so that half of the diffraction peaks vanish. This can be understood as a gradual decrease of the structure factor at these $k$ points.
\begin{figure}[!htbp]
	\centering
	\includegraphics[width=12cm]{blank.png}
\end{figure}

\end{document}
